\chapter{Bonds, Term Structures, and Credit}

\section{Bond cash flows}
A bond is a debt claim. A plain fixed-coupon bond pays coupons and returns principal at maturity. Let $F$ be face value, $c$ the annual coupon rate, $m$ the number of coupon payments per year, $y$ the yield per year compounded at frequency $m$, and $T$ the years to maturity. The coupon per period is $C=Fc/m$ and the number of periods is $N=mT$. Its clean price under the simple model is
\formula{bondprice}
  P=\sum_{k=1}^{N}\frac{C}{(1+y/m)^k}+\frac{F}{(1+y/m)^N}.
\closeformula
Every symbol has a mechanical interpretation: $k$ indexes payment dates, discounting uses the yield convention, and the final principal is discounted to the same date. A zero-coupon bond has no coupons, so $P=F/(1+y/m)^N$.

Price and yield move in opposite directions for a fixed cash-flow schedule. That statement is a property of the discounting formula, not a claim that all observed bonds move together: cash flows, credit quality, liquidity, embedded options, and benchmark conventions also change.

\begin{worked}
An annual coupon bond has face value 1,000, coupon rate 5 percent, and two years to maturity. At a 4 percent yield, $P=50/1.04+1,050/1.04^2\approx1,018.86$. At a 6 percent yield, $P\approx981.67$. The 1 percentage-point yield change causes different price effects depending on the starting yield and cash-flow timing.
\end{worked}

\section{Accrued interest and conventions}
A bond may be quoted on a clean-price basis that excludes accrued coupon, while the buyer pays the dirty price including accrued interest. Day-count conventions determine the fraction of a coupon period that has elapsed. Settlement calendars, business-day adjustments, ex-coupon dates, callable features, and inflation indexation are not cosmetic details; they change cash flows and valuation.

\section{Duration and convexity}
\keyterm{Macaulay duration} is the cash-flow-weighted average time to payment:
\formula{macaulay}
  D_M=\frac{\sum_{k=1}^{N} t_k\,PV(CF_k)}{P},
\closeformula
where $t_k$ is the time in years, $CF_k$ is the cash flow at that date, $PV(CF_k)$ is its present value, and $P$ is bond price. Modified duration under annual compounding is $D_{mod}=D_M/(1+y)$; under $m$ compounding, $D_{mod}=D_M/(1+y/m)$.

For a small yield change $\Delta y$, duration gives the first-order approximation
\formula{durationapprox}
  \frac{\Delta P}{P}\approx-D_{mod}\Delta y.
\closeformula
Convexity captures curvature. Define
\formula{convexity}
  Conv=\frac{1}{P}\sum_{k=1}^{N}\frac{CF_k t_k(t_k+1)}{(1+y)^{t_k+2}}
\closeformula
for annual periods. A second-order approximation is
\formula{convexityapprox}
  \frac{\Delta P}{P}\approx-D_{mod}\Delta y+\frac{1}{2}Conv(\Delta y)^2.
\closeformula
The exact convexity formula changes with frequency and day count. Duration is not a promise of linear price behaviour; it is a local sensitivity.

\section{Spot rates, forward rates, and the curve}
The term structure is the relationship between maturity and interest rates. A spot rate $s_T$ is the discount rate for a single cash flow at time $T$. If a coupon bond's price is known, bootstrapping can infer spot rates from short to long maturities by solving for the discount factor $DF_T$:
\formula{discountfactor}
  P=\sum_{t=1}^{T}CF_tDF_t,\qquad DF_T=\frac{P-\sum_{t=1}^{T-1}CF_tDF_t}{CF_T}.
\closeformula
The implied annual effective spot rate is $s_T=DF_T^{-1/T}-1$. A forward rate is the rate implied between two future dates under the chosen compounding convention. No-arbitrage links spot and forward discount factors, but an implied forward rate is not necessarily an unbiased forecast of the future spot rate.

The curve reflects expected short rates, term premia, inflation risk, liquidity, and supply and demand. The expectations hypothesis is a useful benchmark, not a settled description of every curve movement.

\section{Credit risk}
Credit risk is the possibility that a borrower fails to make promised payments or that its credit quality deteriorates. A simple expected-loss expression is
\formula{expectedloss}
  EL=PD\times LGD\times EAD,
\closeformula
where $PD$ is probability of default over the horizon, $LGD$ is loss given default as a fraction of exposure, and $EAD$ is exposure at default. Recovery may be collateral-dependent, time-dependent, and correlated with default conditions.

Credit spreads compensate for expected loss, risk premia, liquidity, taxes, capital usage, and model uncertainty. A spread is therefore not identical to a default probability. Structural models link default to asset value and leverage; reduced-form models represent default with an intensity or hazard rate. Both are conditional tools.

\section{Inflation-linked and floating-rate bonds}
An inflation-linked bond adjusts principal or coupons using a published index, subject to a lag and legal convention. A floating-rate note resets its coupon against a reference rate plus a spread. Reset frequency, floors, caps, credit spread, and index fallback language matter. A quoted reference rate is not automatically risk-free, and benchmark reform can create contract risk.

\begin{exerciseblock}
\begin{enumerate}
  \item Price a two-year annual-coupon bond with face 1,000, coupon rate 6 percent, and yield 5 percent.
  \item A bond has modified duration 4.2. Estimate its percentage price change when yield rises by 25 basis points.
  \item Explain why convexity improves a duration-only estimate for a larger yield move.
  \item If $PD=2\%$, $LGD=45\%$, and $EAD=5$ million, calculate one-year expected loss before discounting and dependence adjustments.
\end{enumerate}
\end{exerciseblock}

